%% ============================================================
%%  pyWarpFactory — Comprehensive User Manual
%%  LaTeX source
%% ============================================================
\documentclass[11pt, a4paper, twoside]{report}

%% ── Encoding & fonts ────────────────────────────────────────
\usepackage[T1]{fontenc}
\usepackage[utf8]{inputenc}
\usepackage{lmodern}
\usepackage{microtype}

%% ── Page geometry ───────────────────────────────────────────
\usepackage[
  top=2.5cm, bottom=2.8cm,
  inner=3.0cm, outer=2.5cm,
  headheight=14pt
]{geometry}

%% ── Mathematics ─────────────────────────────────────────────
\usepackage{amsmath, amssymb, amsthm, mathtools}
\usepackage{bm}           % bold math
\usepackage{tensor}       % tensor notation  \indices

%% ── Code listings ───────────────────────────────────────────
\usepackage{listings}
\usepackage{xcolor}

\definecolor{codebg}{HTML}{F7F7F7}
\definecolor{codeframe}{HTML}{CCCCCC}
\definecolor{keywordcolor}{HTML}{0000CC}
\definecolor{stringcolor}{HTML}{008000}
\definecolor{commentcolor}{HTML}{808080}
\definecolor{builtincolor}{HTML}{9900CC}

\lstdefinestyle{python}{
  language         = Python,
  backgroundcolor  = \color{codebg},
  frame            = single,
  rulecolor        = \color{codeframe},
  basicstyle       = \ttfamily\small,
  keywordstyle     = \color{keywordcolor}\bfseries,
  stringstyle      = \color{stringcolor},
  commentstyle     = \color{commentcolor}\itshape,
  emph             = {Metric, np, xp, plt, scipy},
  emphstyle        = \color{builtincolor},
  showstringspaces = false,
  breaklines       = true,
  breakatwhitespace= true,
  tabsize          = 4,
  numbers          = left,
  numberstyle      = \tiny\color{gray},
  numbersep        = 6pt,
  xleftmargin      = 1.5em,
  captionpos       = b,
}
\lstset{style=python}

%% ── Coloured boxes ──────────────────────────────────────────
\usepackage{tcolorbox}
\tcbuselibrary{breakable, skins, listings}

\newtcolorbox{infobox}[1][]{
  colback  = blue!5!white,
  colframe = blue!50!black,
  fonttitle= \bfseries,
  title    = Note,
  #1
}
\newtcolorbox{warningbox}[1][]{
  colback  = orange!10!white,
  colframe = orange!70!black,
  fonttitle= \bfseries,
  title    = Warning,
  #1
}
\newtcolorbox{tipbox}[1][]{
  colback  = green!6!white,
  colframe = green!50!black,
  fonttitle= \bfseries,
  title    = Tip,
  #1
}
\newtcolorbox{mathbox}[1][]{
  colback  = gray!7!white,
  colframe = gray!60,
  #1
}

%% ── Tables ──────────────────────────────────────────────────
\usepackage{booktabs, longtable, array, tabularx, multirow}
\renewcommand{\arraystretch}{1.25}

%% ── Graphics & floats ───────────────────────────────────────
\usepackage{graphicx, float}
\usepackage{caption}
\captionsetup{font=small, labelfont=bf}

%% ── Headers & footers ───────────────────────────────────────
\usepackage{fancyhdr}
\pagestyle{fancy}
\fancyhf{}
\fancyhead[LE]{\small\leftmark}
\fancyhead[RO]{\small\rightmark}
\fancyfoot[C]{\small\thepage}
\renewcommand{\headrulewidth}{0.4pt}

%% ── Hyperlinks ──────────────────────────────────────────────
\usepackage[
  colorlinks = true,
  linkcolor  = blue!70!black,
  citecolor  = green!50!black,
  urlcolor   = blue!70!black,
  pdftitle   = {pyWarpFactory User Manual},
  pdfauthor  = {Nelson / xys004},
]{hyperref}
\usepackage{cleveref}

%% ── Miscellaneous ───────────────────────────────────────────
\usepackage{parskip}
\usepackage{enumitem}
\usepackage{titlesec}
\usepackage{pdfpages}   % for possible cover page inclusion

%% ── Custom math macros ──────────────────────────────────────
\newcommand{\dd}{\mathrm{d}}
\newcommand{\pd}[2]{\dfrac{\partial #1}{\partial #2}}
\newcommand{\tensor}[1]{\boldsymbol{#1}}
\newcommand{\Gamma}{\varGamma}
\newcommand{\calL}{\mathcal{L}}
\newcommand{\calM}{\mathcal{M}}
\newcommand{\vect}[1]{\mathbf{#1}}
\newcommand{\diag}{\operatorname{diag}}
\newcommand{\sgn}{\operatorname{sgn}}
\DeclareMathOperator{\tr}{tr}

%% ── Document ────────────────────────────────────────────────
\begin{document}

%% ════════════════════════════════════════════════════════════
%%  TITLE PAGE
%% ════════════════════════════════════════════════════════════
\begin{titlepage}
  \centering
  \vspace*{2cm}

  {\Huge\bfseries pyWarpFactory}\\[0.4em]
  {\Large\itshape A Python Toolkit for Warp Drive Spacetime Analysis}\\[2.5em]

  \rule{0.6\linewidth}{1pt}\\[1.5em]

  {\LARGE\bfseries User Manual}\\[1em]
  {\large Version 1.0 \quad — \quad 2026}\\[3em]

  \begin{minipage}{0.75\textwidth}\centering
    \large
    Numerical framework for generating, solving, and analysing\linebreak
    general-relativistic spacetime metrics, including warp drive\linebreak
    geometries based on the Alcubierre and Van Den Broeck proposals.
  \end{minipage}

  \vfill

  {\large
    \textbf{Author:} Nelson (xys004@gmail.com)\\[0.3em]
    \textbf{Repository:} \url{https://github.com/xys004/pyWarpFactory}\\[0.3em]
    \textbf{License:} MIT
  }

  \vspace{2cm}
\end{titlepage}

\cleardoublepage

%% ════════════════════════════════════════════════════════════
%%  ABSTRACT
%% ════════════════════════════════════════════════════════════
\begin{abstract}
\noindent
\textbf{pyWarpFactory} is a Python port of the MATLAB WarpFactory toolkit
originally developed for the numerical analysis of exotic spacetime metrics
in the context of General Relativity.  The library provides a modular
pipeline covering: (i) metric tensor generation for several physically
motivated spacetimes (Alcubierre warp drive, Schwarzschild black hole,
Van Den Broeck, Lentz soliton, warp shell); (ii) numerical solution of
the Einstein field equations via fourth-order finite differences,
yielding the stress-energy tensor $T^{\mu\nu}$; (iii) Eulerian frame
transformations using tetrad decomposition; and (iv) evaluation of
energy conditions (Null, Weak).  Optional GPU acceleration is supported
through CuPy as a drop-in replacement for NumPy.

This manual provides a complete reference covering installation,
physical and mathematical background, detailed API documentation with
function signatures and return types, worked code examples for every
metric, and practical guidance on grid sizing, unit conventions, and
GPU usage.
\end{abstract}

\cleardoublepage

%% ════════════════════════════════════════════════════════════
%%  TABLE OF CONTENTS
%% ════════════════════════════════════════════════════════════
\tableofcontents
\cleardoublepage

%% ════════════════════════════════════════════════════════════
\chapter{Introduction}
%% ════════════════════════════════════════════════════════════

\section{What is pyWarpFactory?}

pyWarpFactory is a numerical Python library for studying \emph{warp drive}
spacetimes within the framework of General Relativity (GR).  A warp drive
spacetime is a solution to the Einstein field equations that permits
faster-than-light travel by distorting the fabric of spacetime itself,
contracting it in front of the spacecraft and expanding it behind, without
violating local Lorentz invariance.

The toolkit supports the full computational pipeline:
\begin{enumerate}
  \item \textbf{Metric generation} — construct the covariant 4-metric
        $g_{\mu\nu}$ for a chosen spacetime on a discrete grid.
  \item \textbf{Curvature computation} — compute Christoffel symbols
        $\varGamma^\lambda_{\mu\nu}$, Ricci tensor $R_{\mu\nu}$, and
        Ricci scalar $R$ via fourth-order finite differences.
  \item \textbf{Stress-energy tensor} — solve the Einstein field equations
        to obtain $T^{\mu\nu}$, the distribution of energy, momentum, and
        stress implied by the chosen geometry.
  \item \textbf{Energy condition analysis} — evaluate the Null Energy
        Condition (NEC) and Weak Energy Condition (WEC) at every grid
        point to identify regions of exotic matter.
  \item \textbf{ADM decomposition} — extract lapse $\alpha$, shift
        $\beta^i$, and spatial metric $\gamma_{ij}$ from the 3+1 split.
  \item \textbf{Visualisation} — produce publication-quality 2-D slice
        plots of any tensor component.
\end{enumerate}

\section{Historical Background}

Miguel Alcubierre's 1994 paper \emph{``The warp drive: hyper-fast travel
within general relativity''} introduced the first mathematically rigorous
warp drive metric.  Van Den Broeck (1999) subsequently proposed a
modification that dramatically reduces the total exotic energy required.
Lentz (2021) presented a soliton-based solution that, under certain
conditions, claims to avoid exotic matter altogether.  The original
MATLAB WarpFactory was developed to allow systematic numerical exploration
of these metrics.  pyWarpFactory ports that capability to Python,
making it accessible to the broader scientific computing community.

\section{How to Read This Manual}

\begin{itemize}
  \item \textbf{New users} should read \Cref{chap:install,chap:quickstart}
        to get running quickly, then consult \Cref{chap:physics} for the
        mathematical foundations.
  \item \textbf{Intermediate users} interested in a specific metric should
        jump directly to the relevant section in \Cref{chap:generators}.
  \item \textbf{Advanced users} needing full API details should consult
        \Cref{chap:api}.
  \item \Cref{chap:colab} documents the self-contained Google Colab
        notebook included with the project.
\end{itemize}

%% ════════════════════════════════════════════════════════════
\chapter{Installation}
\label{chap:install}
%% ════════════════════════════════════════════════════════════

\section{Requirements}

pyWarpFactory requires Python 3.8 or newer.  The mandatory dependencies
are:

\begin{center}
\begin{tabular}{lll}
  \toprule
  Package & Minimum version & Purpose \\
  \midrule
  \texttt{numpy}      & 1.20.0 & Numerical arrays and linear algebra \\
  \texttt{scipy}      & 1.7.0  & Interpolation, special functions \\
  \texttt{matplotlib} & 3.4.0  & Plotting and visualisation \\
  \bottomrule
\end{tabular}
\end{center}

\textbf{Optional} (GPU acceleration):
\begin{center}
\begin{tabular}{lll}
  \toprule
  Package & Version & Purpose \\
  \midrule
  \texttt{cupy-cuda12x} & $\ge$12 & Drop-in NumPy on CUDA GPU \\
  \bottomrule
\end{tabular}
\end{center}

\section{Installing from Source}

\begin{lstlisting}[caption={Installation from GitHub}]
# Clone the repository
git clone https://github.com/xys004/pyWarpFactory.git
cd pyWarpFactory

# Install (editable mode recommended for development)
pip install -e .

# Or install dependencies only
pip install -r requirements.txt
\end{lstlisting}

\section{Verifying the Installation}

Run the built-in verification script to confirm everything is working
correctly:

\begin{lstlisting}[caption={Package verification}]
python tests/verify_package.py
\end{lstlisting}

This creates a Minkowski flat-space metric, solves for the stress-energy
tensor (which should be identically zero), and reports \texttt{PASS} if
$|T^{\mu\nu}|_\infty < 10^{-10}$.

\section{GPU Setup}

On Google Colab (or any CUDA-enabled machine) CuPy is the recommended
array backend:

\begin{lstlisting}[caption={Install CuPy on Colab (CUDA 12)}]
pip install cupy-cuda12x
\end{lstlisting}

The included \texttt{WarpFactory\_Colab.ipynb} notebook performs this
automatically and falls back to NumPy silently when no GPU is detected.

%% ════════════════════════════════════════════════════════════
\chapter{Quick Start}
\label{chap:quickstart}
%% ════════════════════════════════════════════════════════════

\section{The Simplest Pipeline}

The following five lines run the complete pipeline — from metric creation
to energy condition evaluation — for the Alcubierre warp drive:

\begin{lstlisting}[caption={Minimal Alcubierre example}]
from warpfactory.generator.alcubierre import create_alcubierre_metric
from warpfactory.analyzer.energy_conditions import (
    calculate_energy_conditions
)

# 1. Grid definition
grid_size  = (7, 30, 30, 30)        # (Nt, Nx, Ny, Nz)
grid_scale = (1.0, 0.3, 0.3, 0.3)  # (dt, dx, dy, dz) code units

# 2. World centre (physical coords of grid centre)
world_center = (1.0, 4.5, 4.5, 4.5)

# 3. Generate metric
metric = create_alcubierre_metric(
    grid_size, grid_scale, world_center,
    v=0.5,   # bubble speed (fraction of C)
    R=1.5,   # bubble radius
    sigma=8.0
)

# 4. Solve + evaluate energy conditions
results, T_euler = calculate_energy_conditions(
    metric.tensor, grid_scale
)

# 5. Inspect
import numpy as np
print("Min NEC:", np.min(results['Null']))   # negative = violated
print("Energy density min:", np.min(T_euler[0, 0]))
\end{lstlisting}

\section{Visualising Results}

\begin{lstlisting}[caption={Plotting a 2-D slice of the energy density}]
import matplotlib.pyplot as plt

t_sl, z_sl = 3, 15               # time and z slice indices
rho = T_euler[0, 0, t_sl, :, :, z_sl]  # shape (Nx, Ny)

plt.figure(figsize=(6, 5))
plt.imshow(rho.T, origin='lower', cmap='RdBu_r',
           extent=[0, 30*0.3, 0, 30*0.3])
plt.colorbar(label=r'$\rho = \hat{T}^{00}$')
plt.xlabel('x')
plt.ylabel('y')
plt.title('Energy Density — Alcubierre Warp Drive')
plt.tight_layout()
plt.savefig('energy_density.png', dpi=130)
plt.show()
\end{lstlisting}

\section{Understanding the Output}

\paragraph{The \texttt{results} dictionary} contains two keys:
\begin{itemize}
  \item \texttt{results['Null']}: the minimum value of
        $T_{\mu\nu}k^\mu k^\nu$ over a dense sample of null vectors $k$,
        evaluated at every grid point.  \textbf{Negative values indicate
        a Null Energy Condition (NEC) violation.}
  \item \texttt{results['Weak']}: $\min(\rho,\;\text{NEC})$, combining
        the energy density requirement with the NEC.
        \textbf{Negative values indicate a Weak Energy Condition (WEC)
        violation.}
\end{itemize}

\paragraph{The \texttt{T\_euler} array} has shape
\texttt{(4, 4, Nt, Nx, Ny, Nz)} and stores the contravariant
stress-energy tensor $\hat{T}^{\hat\mu\hat\nu}$ in the Eulerian (locally
flat) frame.  The energy density seen by a static local observer is
$\rho = \hat{T}^{00}$.

%% ════════════════════════════════════════════════════════════
\chapter{Physical and Mathematical Background}
\label{chap:physics}
%% ════════════════════════════════════════════════════════════

\section{Conventions}

Throughout this manual and in the code, we use the following conventions:

\begin{itemize}
  \item \textbf{Signature:} $(-,+,+,+)$ (mostly-positive, or ``East
        Coast'') — the metric is $g_{\mu\nu} = \diag(-1,1,1,1)$ in flat
        space.
  \item \textbf{Index order:} $\mu = 0$ is the time component;
        $\mu=1,2,3$ are $x,y,z$ respectively.
  \item \textbf{Units:} Code units $c=G=1$ are used throughout unless
        explicitly stated.  Physical SI values are stored in
        \texttt{warpfactory/constants.py}.
  \item \textbf{Einstein summation:} repeated upper/lower indices are
        summed over $0,1,2,3$.
\end{itemize}

\section{The Einstein Field Equations}

The Einstein field equations relate spacetime curvature to energy-momentum:
\begin{mathbox}
\begin{equation}
  G_{\mu\nu} \;=\; R_{\mu\nu} - \tfrac{1}{2}\,R\,g_{\mu\nu}
             \;=\; \frac{8\pi G}{c^4}\,T_{\mu\nu}
  \label{eq:EFE}
\end{equation}
\end{mathbox}
where $G_{\mu\nu}$ is the Einstein tensor, $R_{\mu\nu}$ the Ricci tensor,
$R$ the Ricci scalar, and $T_{\mu\nu}$ the stress-energy tensor.

\subsection{Christoffel Symbols}

Given the metric $g_{\mu\nu}$, the Levi-Civita connection coefficients
(Christoffel symbols of the second kind) are:
\begin{equation}
  \varGamma^\lambda_{\mu\nu}
  = \tfrac{1}{2}\,g^{\lambda\rho}
    \bigl(\partial_\mu g_{\nu\rho}
         + \partial_\nu g_{\mu\rho}
         - \partial_\rho g_{\mu\nu}\bigr)
  \label{eq:christoffel}
\end{equation}

\subsection{Ricci Tensor}

\begin{equation}
  R_{\mu\nu}
  = \partial_\rho \varGamma^\rho_{\mu\nu}
  - \partial_\nu  \varGamma^\rho_{\mu\rho}
  + \varGamma^\rho_{\mu\nu}\varGamma^\sigma_{\rho\sigma}
  - \varGamma^\sigma_{\mu\rho}\varGamma^\rho_{\nu\sigma}
  \label{eq:ricci}
\end{equation}

\subsection{Ricci Scalar}

\begin{equation}
  R = g^{\mu\nu} R_{\mu\nu}
  \label{eq:riccisc}
\end{equation}

\subsection{Stress-Energy Tensor}

Inverting \cref{eq:EFE}:
\begin{equation}
  T_{\mu\nu} = \frac{c^4}{8\pi G}\,G_{\mu\nu},
  \qquad
  T^{\mu\nu} = g^{\mu\alpha}g^{\nu\beta}T_{\alpha\beta}
  \label{eq:SET}
\end{equation}

\section{Numerical Differentiation}

All partial derivatives are computed using the \emph{fourth-order central
finite difference} stencil:
\begin{equation}
  \pd{f}{x}\bigg|_i
  \;\approx\;
  \frac{-f_{i+2} + 8f_{i+1} - 8f_{i-1} + f_{i-2}}{12\,\Delta x}
  + \mathcal{O}(\Delta x^4)
  \label{eq:fd4}
\end{equation}
This requires at least five grid points in the differentiation direction.
Boundary cells (indices 0, 1 and $N{-}2$, $N{-}1$) are filled with the
nearest valid interior derivative value (first-order accurate at the
boundary), which suffices for the smooth metric profiles typically used.

\section{The 3+1 ADM Decomposition}
\label{sec:adm}

The Arnowitt--Deser--Misner (ADM) decomposition splits the spacetime
metric into spatial and temporal parts via:
\begin{equation}
  \dd s^2
  = -\alpha^2\,\dd t^2
  + \gamma_{ij}\!\left(\dd x^i + \beta^i\,\dd t\right)\!
                     \left(\dd x^j + \beta^j\,\dd t\right)
  \label{eq:adm}
\end{equation}
where:
\begin{itemize}
  \item $\alpha$ is the \textbf{lapse function} — measures proper time
        elapsed per unit coordinate time.
  \item $\beta^i$ is the \textbf{shift vector} — describes how spatial
        coordinates shift between successive time slices.
  \item $\gamma_{ij}$ is the \textbf{spatial metric} induced on constant-$t$
        hypersurfaces.
\end{itemize}

The component relations are:
\begin{align}
  g_{00} &= -\alpha^2 + \gamma_{ij}\beta^i\beta^j, &
  g_{0i} &= \gamma_{ij}\beta^j = \beta_i, &
  g_{ij} &= \gamma_{ij}
  \label{eq:adm_comp}
\end{align}

\section{Energy Conditions}
\label{sec:ec}

Energy conditions are constraints on $T_{\mu\nu}$ that classical forms
of matter are expected to satisfy.  Warp drive metrics typically
\emph{violate} them, implying the need for exotic matter.

\subsection{Null Energy Condition (NEC)}
For every future-directed null vector $k^\mu$ ($k^\mu k_\mu = 0$):
\begin{equation}
  T_{\mu\nu}\,k^\mu k^\nu \;\ge\; 0
  \label{eq:nec}
\end{equation}

\subsection{Weak Energy Condition (WEC)}
For every future-directed timelike vector $t^\mu$ ($t^\mu t_\mu < 0$):
\begin{equation}
  T_{\mu\nu}\,t^\mu t^\nu \;\ge\; 0
  \label{eq:wec}
\end{equation}
In the Eulerian frame this reduces to $\rho \ge 0$ combined with the NEC
(the null direction is the limiting timelike case).

\subsection{Eulerian Observer}
The Eulerian (normal) observer is at rest in the spatial coordinates.
Its 4-velocity is $n^\mu = (1/\alpha, -\beta^i/\alpha)$, and the energy
density it measures is:
\begin{equation}
  \rho = T_{\mu\nu}\,n^\mu n^\nu
       = \hat{T}^{00}
  \label{eq:rho_eulerian}
\end{equation}
where $\hat{T}^{\hat\mu\hat\nu}$ denotes components in the locally flat
Eulerian tetrad frame.

\section{Eulerian Frame Transformation}
\label{sec:euleriantf}

To transform the stress-energy tensor to the Eulerian tetrad frame,
pyWarpFactory uses a Cholesky-like decomposition.  The transformation
matrix $M^a{}_\mu$ satisfies $M^a{}_\mu M^b{}_\nu g^{\mu\nu} = \eta^{ab}$
where $\eta = \diag(-1,+1,+1,+1)$.  The Eulerian frame components are:
\begin{equation}
  \hat{T}_{ab} = M^{\mu}{}_{a}\,T_{\mu\nu}\,M^{\nu}{}_{b}
\end{equation}
The explicit matrix $M$ is computed column-by-column using the recursive
formulas derived from the metric determinant factorisation (see
\texttt{get\_eulerian\_transformation\_matrix} in \S\ref{sec:api_transform}).

%% ════════════════════════════════════════════════════════════
\chapter{Package Structure}
\label{chap:structure}
%% ════════════════════════════════════════════════════════════

\section{Directory Layout}

\begin{lstlisting}[language=bash, caption={Repository layout}]
pyWarpFactory/
|-- warpfactory/              # Main Python package
|   |-- __init__.py
|   |-- constants.py          # Physical constants (C, G)
|   |-- generator/            # Metric constructors
|   |   |-- base.py           # Metric dataclass
|   |   |-- commons.py        # Shared utilities / shape functions
|   |   |-- alcubierre.py
|   |   |-- schwarzschild.py
|   |   |-- van_den_broeck.py
|   |   |-- lentz.py
|   |   |-- warp_shell.py
|   |   |-- minkowski.py
|   |   `-- modified_time.py
|   |-- solver/               # Einstein equation solvers
|   |   |-- solvers.py        # High-level solve_energy_tensor()
|   |   |-- tensor_utils.py   # Christoffel, Ricci, Einstein, ...
|   |   `-- finite_difference.py
|   |-- analyzer/             # Physical analysis
|   |   |-- energy_conditions.py
|   |   |-- transform.py      # Frame transformations
|   |   |-- scalars.py        # ADM scalars (expansion, shear, ...)
|   |   `-- momentum_flow.py
|   |-- algo/
|   |   `-- adm.py            # ADM decompose / reconstruct
|   |-- utils/
|   |   `-- helpers.py        # TOV, interpolation, Legendre, ...
|   `-- visualizer/
|       `-- plotter.py
|-- tests/                    # Validation scripts
|-- WarpFactory_Colab.ipynb   # Self-contained Google Colab notebook
|-- setup.py
|-- requirements.txt
`-- USER_MANUAL_LATEX.tex     # This document
\end{lstlisting}

\section{Core Data Structure: the \texttt{Metric} Class}
\label{sec:metric_class}

Every object produced by the generator module (and most objects returned
by the solver and analyser) is an instance of the \texttt{Metric} dataclass
defined in \texttt{warpfactory/generator/base.py}.

\begin{lstlisting}[caption={Metric dataclass definition (simplified)}]
from dataclasses import dataclass, field
from typing import Dict, Any
import numpy as np

@dataclass
class Metric:
    tensor  : np.ndarray          # shape (4, 4, Nt, Nx, Ny, Nz)
    coords  : Dict[str, np.ndarray]  # {'t':…, 'x':…, 'y':…, 'z':…}
    scaling : np.ndarray          # [dt, dx, dy, dz]
    name    : str  = "Generic"
    index   : str  = "covariant"  # or "contravariant"
    params  : Dict[str, Any] = field(default_factory=dict)
\end{lstlisting}

\paragraph{Key attributes:}
\begin{description}
  \item[\texttt{tensor}]
    The 4-metric $g_{\mu\nu}$ (or any rank-2 tensor) stored as a
    NumPy/CuPy array of shape \texttt{(4, 4, Nt, Nx, Ny, Nz)}.
    Axis 0 is the first spacetime index $\mu$; axis 1 is $\nu$;
    axes 2--5 are the grid dimensions $(t, x, y, z)$.

  \item[\texttt{coords}]
    Dictionary with keys \texttt{'t'}, \texttt{'x'}, \texttt{'y'},
    \texttt{'z'}, each an array of shape \texttt{(Nt, Nx, Ny, Nz)}
    storing the physical (centred) coordinate at every grid point.

  \item[\texttt{scaling}]
    1-D array \texttt{[dt, dx, dy, dz]} — the uniform grid spacings
    used by the finite-difference solver.

  \item[\texttt{index}]
    Either \texttt{'covariant'} (subscript indices, $g_{\mu\nu}$)
    or \texttt{'contravariant'} (superscript, $T^{\mu\nu}$).

  \item[\texttt{params}]
    Metric-specific parameters, e.g.\
    $\{v, R, \sigma\}$ for Alcubierre.
\end{description}

\section{Grid Conventions}

\begin{itemize}
  \item Coordinates start at zero: \texttt{x[i] = i * dx}.
  \item A 1-cell offset is applied inside the generator functions to
        match MATLAB's 1-based indexing convention:
        \texttt{x\_phys = (x + dx) - x\_center}.
  \item \textbf{Minimum grid sizes:} at least 5 points in every
        direction that is differentiated (all four for the solver).
        Typical choices: $N_t = 7$, $N_x = N_y = N_z \ge 25$.
\end{itemize}

%% ════════════════════════════════════════════════════════════
\chapter{Metric Generators}
\label{chap:generators}
%% ════════════════════════════════════════════════════════════

\section{Common Infrastructure}

\subsection{\texttt{create\_grid}}
\begin{lstlisting}
from warpfactory.generator.commons import create_grid
grid = create_grid(grid_size=(Nt, Nx, Ny, Nz),
                   grid_scale=(dt, dx, dy, dz))
# Returns: dict {'t': array, 'x': array, 'y': array, 'z': array}
# Each array has shape (Nt, Nx, Ny, Nz) using 'ij' indexing.
\end{lstlisting}

\subsection{\texttt{shape\_function\_alcubierre}}
The Alcubierre shape function $f(r; R, \sigma)$ defines the bubble profile:
\begin{equation}
  f(r;\,R,\sigma)
  = \frac{\tanh\!\bigl[\sigma(R+r)\bigr]
         + \tanh\!\bigl[\sigma(R-r)\bigr]}
         {2\,\tanh(\sigma R)}
  \label{eq:shape_fn}
\end{equation}
$f \approx 1$ inside radius $R$ and decays smoothly to $0$ outside;
$\sigma$ controls the shell thickness.

\begin{lstlisting}
from warpfactory.generator.commons import shape_function_alcubierre
import numpy as np
r = np.linspace(0, 5, 300)
f = shape_function_alcubierre(r, R=2.0, sigma=8.0)
\end{lstlisting}

%% ────────────────────────────────────────────────────────────
\section{Minkowski (Flat Spacetime)}
%% ────────────────────────────────────────────────────────────

\subsection{Physics}

Flat Minkowski spacetime is the trivial vacuum solution:
\begin{equation}
  g_{\mu\nu} = \eta_{\mu\nu} = \diag(-1,+1,+1,+1)
\end{equation}
The stress-energy tensor is identically zero everywhere.

\subsection{API}

\begin{lstlisting}[caption={Minkowski metric}]
from warpfactory.generator.minkowski import create_minkowski_metric

metric = create_minkowski_metric(
    grid_size  = (7, 20, 20, 20),
    grid_scale = (1.0, 0.2, 0.2, 0.2)
)
\end{lstlisting}

\begin{description}
  \item[\texttt{grid\_size}] \texttt{(Nt, Nx, Ny, Nz)} — number of
        grid points in each dimension.
  \item[\texttt{grid\_scale}] \texttt{(dt, dx, dy, dz)} — spacing in
        code units.
  \item[Returns] \texttt{Metric}, \texttt{index='covariant'},
        \texttt{name='Minkowski'}.
\end{description}

Use this metric as a sanity check: after running the full solver
pipeline, all energy condition values should be $\approx 0$.

%% ────────────────────────────────────────────────────────────
\section{Alcubierre Warp Drive}
\label{sec:alcubierre}
%% ────────────────────────────────────────────────────────────

\subsection{Physics}

The Alcubierre metric (1994) describes a bubble of flat spacetime moving
at arbitrary speed $v$ in the $x$-direction.  In ADM form with lapse
$\alpha=1$, flat spatial metric $\gamma_{ij}=\delta_{ij}$, and shift:
\begin{equation}
  \beta^x = -v\,f(r_s),
  \qquad \beta^y = \beta^z = 0,
\end{equation}
where $r_s = \sqrt{(x - x_s(t))^2 + y^2 + z^2}$ is the distance from
the moving bubble centre $x_s(t) = v\,t$.  The covariant 4-metric is:
\begin{align}
  g_{00} &= -1 + (v\,f)^2, &
  g_{01} = g_{10} &= -v\,f, \label{eq:alc_metric} \\
  g_{11} = g_{22} = g_{33} &= 1, &
  g_{\mu\neq\nu, \,\mu,\nu\ge1} &= 0. \nonumber
\end{align}

Observers inside the bubble travel effectively at $v > c$ without
experiencing superluminal motion locally.  The bubble shell requires
\emph{negative energy density} (exotic matter), violating both the NEC
and WEC.

\subsection{API}

\begin{lstlisting}[caption={Alcubierre metric}]
from warpfactory.generator.alcubierre import create_alcubierre_metric

metric = create_alcubierre_metric(
    grid_size    = (7, 30, 30, 30),
    grid_scale   = (1.0, 0.3, 0.3, 0.3),
    world_center = (1.0, 4.5, 4.5, 4.5),  # (tc, xc, yc, zc)
    v     = 0.5,    # bubble speed (in units of C)
    R     = 1.5,    # bubble radius
    sigma = 8.0     # shell steepness
)
\end{lstlisting}

\begin{center}
\begin{tabularx}{\linewidth}{llX}
  \toprule
  Parameter & Type & Description \\
  \midrule
  \texttt{grid\_size}    & tuple(4) & $(N_t, N_x, N_y, N_z)$ grid dimensions \\
  \texttt{grid\_scale}   & tuple(4) & $(dt, dx, dy, dz)$ grid spacing \\
  \texttt{world\_center} & tuple(4) & $(t_c, x_c, y_c, z_c)$ physical centre \\
  \texttt{v}             & float    & Bubble velocity in units of $C$ \\
  \texttt{R}             & float    & Bubble radius (same units as grid) \\
  \texttt{sigma}         & float    & Shell steepness; larger $= $ thinner shell \\
  \midrule
  \emph{Returns}         & Metric   & Covariant $g_{\mu\nu}$, \texttt{name='Alcubierre'} \\
  \bottomrule
\end{tabularx}
\end{center}

\begin{warningbox}
When using SI units ($C = 2.998\times10^8$ m/s) the bubble travels at
$v \cdot C$ metres per second.  For typical grid sizes of a few metres,
the bubble exits the domain in nanoseconds.  Use code units ($C=1$) for
all exploratory work and rescale only when comparing to SI results.
\end{warningbox}

%% ────────────────────────────────────────────────────────────
\section{Schwarzschild Black Hole}
\label{sec:schwarzschild}
%% ────────────────────────────────────────────────────────────

\subsection{Physics}

The Schwarzschild metric describes the unique spherically symmetric
vacuum solution with Schwarzschild radius $r_s = 2GM/c^2$:
\begin{equation}
  \dd s^2
  = -\!\left(1 - \frac{r_s}{r}\right)\dd t^2
  + \left(1 - \frac{r_s}{r}\right)^{\!-1}\!\dd r^2
  + r^2\,\dd\Omega^2
  \label{eq:schw_spherical}
\end{equation}
pyWarpFactory stores this in quasi-Cartesian coordinates $(t,x,y,z)$
with off-diagonal spatial terms to encode the angular part.  The
stress-energy tensor is identically zero everywhere outside the horizon
(pure vacuum), so energy conditions are trivially satisfied in the
exterior region.

\subsection{API}

\begin{lstlisting}[caption={Schwarzschild metric}]
from warpfactory.generator.schwarzschild import create_schwarzschild_metric

# Note: Nt must be 1 (static metric)
metric = create_schwarzschild_metric(
    grid_size    = (1, 25, 25, 25),
    grid_scale   = (1.0, 0.3, 0.3, 0.3),
    world_center = (1.0, 3.75, 3.75, 3.75),
    rs           = 1.0   # Schwarzschild radius (code units: rs = 2M)
)
\end{lstlisting}

\begin{center}
\begin{tabularx}{\linewidth}{llX}
  \toprule
  Parameter & Type & Description \\
  \midrule
  \texttt{grid\_size}    & tuple(4) & Must have $N_t = 1$ (static) \\
  \texttt{grid\_scale}   & tuple(4) & $(dt, dx, dy, dz)$ \\
  \texttt{world\_center} & tuple(4) & Spatial origin of the black hole \\
  \texttt{rs}            & float    & Schwarzschild radius; code units: $r_s = 2M$ \\
  \midrule
  \emph{Returns}         & Metric   & Covariant metric, \texttt{name='Schwarzschild'} \\
  \bottomrule
\end{tabularx}
\end{center}

\begin{infobox}
An \texttt{epsilon = 1e-10} guard prevents division by zero at the
coordinate singularity $r = 0$.  Keep the singularity away from the grid
boundary by choosing an appropriate \texttt{world\_center}.
\end{infobox}

%% ────────────────────────────────────────────────────────────
\section{Van Den Broeck Metric}
\label{sec:vdb}
%% ────────────────────────────────────────────────────────────

\subsection{Physics}

Van Den Broeck (1999) modified the Alcubierre metric by expanding
the spatial volume inside the bubble by a large factor $(1+A)^2$,
making the interior macroscopically large while keeping the exterior
bubble macroscopically small.  This dramatically reduces the total
exotic energy required.

The spatial metric components are $\gamma_{ij} = B^2(r)\delta_{ij}$
where:
\begin{equation}
  B(r) = 1 + A\,f(r;\,R_1,\sigma_1)
\end{equation}
The shift vector uses the standard Alcubierre form but with an
effective velocity $v_\text{eff} = v(1+A)^2$:
\begin{equation}
  g_{00} = -\!\left[1 - B^2 f_s^2\right],\quad
  g_{01} = -B^2 f_s, \quad
  g_{11} = g_{22} = g_{33} = B^2
\end{equation}
where $f_s = f(r;\,R_2,\sigma_2)\cdot v$.

\subsection{API}

\begin{lstlisting}[caption={Van Den Broeck metric}]
from warpfactory.generator.van_den_broeck import (
    create_van_den_broeck_metric
)

metric = create_van_den_broeck_metric(
    grid_size    = (7, 30, 30, 30),
    grid_scale   = (1.0, 0.3, 0.3, 0.3),
    world_center = (1.0, 4.5, 4.5, 4.5),
    v      = 0.5,   # base velocity
    R1     = 2.5,   # spatial expansion radius
    sigma1 = 6.0,   # expansion steepness
    R2     = 1.0,   # shift-vector radius
    sigma2 = 8.0,   # shift steepness
    A      = 2.0    # spatial expansion factor
)
\end{lstlisting}

\begin{center}
\begin{tabularx}{\linewidth}{llX}
  \toprule
  Parameter & Type & Description \\
  \midrule
  \texttt{v}      & float & Base velocity \\
  \texttt{R1, sigma1} & float & Expansion region shape parameters \\
  \texttt{R2, sigma2} & float & Shift-vector region shape parameters \\
  \texttt{A}      & float & Spatial expansion factor; effective speed
                            $v_\text{eff} = v(1+A)^2$ \\
  \midrule
  \emph{Returns}  & Metric & Covariant, \texttt{name='Van Den Broeck'} \\
  \bottomrule
\end{tabularx}
\end{center}

%% ────────────────────────────────────────────────────────────
\section{Lentz Soliton}
\label{sec:lentz}
%% ────────────────────────────────────────────────────────────

\subsection{Physics}

Lentz (2021) proposed a shift-vector soliton built from piecewise
regions with alternating positive and negative warp factors.  The
structure is defined on a 2-D template with seven distinct regions
using values in $\{-2,-1,0,0.5,1\}$.  Under certain conditions it
was claimed to avoid requiring exotic energy.

\subsection{API}

\begin{lstlisting}[caption={Lentz soliton metric}]
from warpfactory.generator.lentz import create_lentz_metric

metric = create_lentz_metric(
    grid_size    = (7, 30, 30, 30),
    grid_scale   = (1.0, 0.3, 0.3, 0.3),
    world_center = (1.0, 4.5, 4.5, 4.5),
    v     = 0.5,
    scale = None    # auto: max(Nx,Ny,Nz)/7
)
\end{lstlisting}

\begin{center}
\begin{tabularx}{\linewidth}{llX}
  \toprule
  Parameter & Type & Description \\
  \midrule
  \texttt{v}     & float       & Soliton velocity \\
  \texttt{scale} & float/None  & Size factor; default = $\max(N)/7$ \\
  \midrule
  \emph{Returns} & Metric      & Covariant, \texttt{name='Lentz'} \\
  \bottomrule
\end{tabularx}
\end{center}

%% ────────────────────────────────────────────────────────────
\section{Modified Time Metric}
\label{sec:modtime}
%% ────────────────────────────────────────────────────────────

\subsection{Physics}

A variant of the Alcubierre metric that modifies the lapse function
inside the bubble using an additional parameter $A$:
\begin{equation}
  g_{00} = -\left[(1 - f) + \tfrac{f}{A}\right]^2 + (v\,f)^2, \quad
  g_{01} = -v\,f, \quad g_{11} = g_{22} = g_{33} = 1
\end{equation}
When $A=1$ this reduces to the standard Alcubierre metric.

\subsection{API}

\begin{lstlisting}[caption={Modified time metric}]
from warpfactory.generator.modified_time import (
    create_modified_time_metric
)

metric = create_modified_time_metric(
    grid_size    = (7, 25, 25, 25),
    grid_scale   = (1.0, 0.3, 0.3, 0.3),
    world_center = (1.0, 3.75, 3.75, 3.75),
    v     = 0.5,
    R     = 1.5,
    sigma = 8.0,
    A     = 1.5    # lapse modification factor
)
\end{lstlisting}

%% ────────────────────────────────────────────────────────────
\section{Warp Shell}
\label{sec:warpshell}
%% ────────────────────────────────────────────────────────────

\subsection{Physics}

The warp shell metric models a massive spherical shell of matter using
the Tolman-Oppenheimer-Volkoff (TOV) equation for a constant-density
sphere.  It serves as a physically realistic background for testing the
solver and energy condition machinery.

The metric is generated in spherical coordinates and transformed to
Cartesian $(x, y, z)$ using the transformation described in
\texttt{helpers.sph2cart\_diag}.  Optionally a warp effect
(\texttt{do\_warp=True}) can be superimposed.

\subsection{API}

\begin{lstlisting}[caption={Warp shell metric}]
from warpfactory.generator.warp_shell import create_warp_shell_metric

metric = create_warp_shell_metric(
    grid_size    = (7, 25, 25, 25),
    grid_scale   = (1.0, 0.3, 0.3, 0.3),
    world_center = (1.0, 3.75, 3.75, 3.75),
    mass         = 1.0,     # total shell mass
    r_inner      = 0.5,     # inner radius
    r_outer      = 1.5,     # outer radius
    r_buff       = 0.2,     # buffer distance
    sigma        = 10.0,    # sigmoid sharpness
    smooth_factor= 2,       # smoothing iterations
    v_warp       = 0.0,     # warp velocity (if do_warp=True)
    do_warp      = False    # enable warp effect
)
\end{lstlisting}

\begin{center}
\begin{tabularx}{\linewidth}{llX}
  \toprule
  Parameter      & Type  & Description \\
  \midrule
  \texttt{mass}         & float & Total shell mass ($M$) \\
  \texttt{r\_inner}     & float & Inner radius of the shell \\
  \texttt{r\_outer}     & float & Outer radius of the shell \\
  \texttt{r\_buff}      & float & Transition buffer width \\
  \texttt{sigma}        & float & Sigmoid sharpness at interfaces \\
  \texttt{smooth\_factor} & int & Number of Legendre smoothing passes \\
  \texttt{v\_warp}      & float & Warp velocity (if \texttt{do\_warp=True}) \\
  \texttt{do\_warp}     & bool  & Whether to apply a warp shift \\
  \midrule
  \emph{Returns} & Metric & Covariant, \texttt{name='Warp Shell'} \\
  \bottomrule
\end{tabularx}
\end{center}

%% ════════════════════════════════════════════════════════════
\chapter{The Solver Module}
\label{chap:solver}
%% ════════════════════════════════════════════════════════════

\section{Overview}

The solver module computes the stress-energy tensor $T^{\mu\nu}$ from a
given metric by numerically evaluating the Einstein field equations
(\cref{eq:EFE}).  The computational pipeline is:

\begin{center}
\begin{tabular}{ccccc}
  $g_{\mu\nu}$ & $\to$ & $\varGamma^\lambda_{\mu\nu}$ & $\to$ &
  $R_{\mu\nu}$ \\[0.3em]
  & & & $\to$ & $R$ \\[0.3em]
  & & & $\to$ & $G_{\mu\nu}$ \\[0.3em]
  & & & $\to$ & $T^{\mu\nu}$
\end{tabular}
\end{center}

\section{High-Level Function: \texttt{solve\_energy\_tensor}}

\begin{lstlisting}[caption={Full Einstein pipeline}]
from warpfactory.solver.solvers import solve_energy_tensor

energy_tensor = solve_energy_tensor(metric)
# Returns: Metric object with type='Stress-Energy', index='contravariant'
# energy_tensor.tensor has shape (4, 4, Nt, Nx, Ny, Nz)
\end{lstlisting}

\begin{description}
  \item[Input] A \texttt{Metric} object (any generator output).
  \item[Output] A new \texttt{Metric} object containing $T^{\mu\nu}$
        with \texttt{index='contravariant'} and
        \texttt{type='Stress-Energy'}.
\end{description}

\begin{infobox}
\texttt{solve\_energy\_tensor} is the preferred high-level interface.
The lower-level functions \texttt{get\_christoffel}, \texttt{get\_ricci\_tensor},
etc.\ are available for custom pipelines.
\end{infobox}

\section{Low-Level Functions}

\subsection{\texttt{get\_christoffel}}
\begin{lstlisting}
from warpfactory.solver.tensor_utils import get_christoffel

gamma = get_christoffel(metric.tensor, metric.scaling)
# gamma.shape == (4, 4, 4, Nt, Nx, Ny, Nz)
# Axes: [lambda, mu, nu, t, x, y, z]
\end{lstlisting}
Computes $\varGamma^\lambda_{\mu\nu}$ using \cref{eq:christoffel} with
fourth-order finite differences.

\subsection{\texttt{get\_ricci\_tensor}}
\begin{lstlisting}
from warpfactory.solver.tensor_utils import get_ricci_tensor

R_mn = get_ricci_tensor(metric.tensor, metric.scaling)
# R_mn.shape == (4, 4, Nt, Nx, Ny, Nz)
\end{lstlisting}

\subsection{\texttt{get\_ricci\_scalar}}
\begin{lstlisting}
from warpfactory.solver.tensor_utils import (
    get_ricci_scalar, get_c4_inv
)
g_inv = get_c4_inv(metric.tensor)
R = get_ricci_scalar(R_mn, g_inv)
# R.shape == (Nt, Nx, Ny, Nz)
\end{lstlisting}

\subsection{\texttt{get\_einstein\_tensor}}
\begin{lstlisting}
from warpfactory.solver.tensor_utils import get_einstein_tensor

G_mn = get_einstein_tensor(R_mn, R, metric.tensor)
# G_mn.shape == (4, 4, Nt, Nx, Ny, Nz)
\end{lstlisting}

\subsection{\texttt{get\_energy\_tensor}}
\begin{lstlisting}
from warpfactory.solver.tensor_utils import get_energy_tensor

T_uv = get_energy_tensor(G_mn, g_inv)
# T_uv.shape == (4, 4, Nt, Nx, Ny, Nz)  -- contravariant T^{mu nu}
\end{lstlisting}

\subsection{\texttt{get\_c4\_inv} and \texttt{get\_c3\_inv}}
\begin{lstlisting}
from warpfactory.solver.tensor_utils import get_c4_inv, get_c3_inv

g_inv      = get_c4_inv(metric.tensor)          # (4,4,...) -> (4,4,...)
gamma_inv  = get_c3_inv(spatial_metric)          # (3,3,...) -> (3,3,...)
\end{lstlisting}
Both functions invert the tensor pointwise using
\texttt{numpy.linalg.inv} with axis reordering.

\section{Finite Difference Module}

\subsection{\texttt{derive\_1st\_4th\_order}}
\begin{lstlisting}
from warpfactory.solver.finite_difference import derive_1st_4th_order

dA_dx = derive_1st_4th_order(A, axis=2, delta=dx)
\end{lstlisting}
Computes $\partial A / \partial x^{\text{axis}}$ using the stencil of
\cref{eq:fd4}.  Returns \texttt{zeros\_like(A)} if
\texttt{A.shape[axis] < 5}.

\subsection{\texttt{derive\_2nd\_4th\_order}}
\begin{lstlisting}
from warpfactory.solver.finite_difference import derive_2nd_4th_order

# Unmixed: d²A/dx²
d2A = derive_2nd_4th_order(A, axis1=2, axis2=2, delta1=dx)

# Mixed: d²A/(dx dy)
d2A = derive_2nd_4th_order(A, axis1=2, axis2=3, delta1=dx, delta2=dy)
\end{lstlisting}

\begin{center}
\begin{tabularx}{\linewidth}{llX}
  \toprule
  Parameter    & Type  & Description \\
  \midrule
  \texttt{A}       & ndarray & Array to differentiate \\
  \texttt{axis}    & int     & Axis index (0-based); note axes 0,1 are
                               tensor indices, axis 2 = $t$, 3 = $x$, \ldots \\
  \texttt{delta}   & float   & Grid spacing along that axis \\
  \midrule
  \emph{Returns} & ndarray & Same shape as \texttt{A} \\
  \bottomrule
\end{tabularx}
\end{center}

\begin{warningbox}
Axes 0 and 1 in the \texttt{Metric.tensor} array are the tensor component
indices $(\mu, \nu)$.  Physical coordinates start at axis 2 ($t$), axis 3
($x$), axis 4 ($y$), axis 5 ($z$).  Always pass \texttt{axis = \{2,3,4,5\}}
when differentiating metric components.
\end{warningbox}

%% ════════════════════════════════════════════════════════════
\chapter{The Analyzer Module}
\label{chap:analyzer}
%% ════════════════════════════════════════════════════════════

\section{Energy Conditions}
\label{sec:ec_api}

\subsection{Main Function: \texttt{calculate\_energy\_conditions}}

This is the primary high-level function.  It runs the full pipeline
from a raw metric tensor to condition violation maps.

\begin{lstlisting}[caption={Energy conditions — full pipeline}]
from warpfactory.analyzer.energy_conditions import (
    calculate_energy_conditions
)

results, T_euler = calculate_energy_conditions(
    metric_tensor = metric.tensor,     # (4,4,Nt,Nx,Ny,Nz)
    grid_scale    = metric.scaling,    # (dt,dx,dy,dz)
    num_vecs      = 50                 # null vector sample size
)
\end{lstlisting}

\begin{center}
\begin{tabularx}{\linewidth}{llX}
  \toprule
  Parameter        & Type    & Description \\
  \midrule
  \texttt{metric\_tensor} & ndarray & Covariant 4-metric, shape
                                    \texttt{(4,4,Nt,Nx,Ny,Nz)} \\
  \texttt{grid\_scale}    & tuple   & $(dt, dx, dy, dz)$ \\
  \texttt{num\_vecs}      & int     & Number of null test vectors on the
                                    sphere (default 50; more = tighter
                                    lower bound) \\
  \midrule
  \texttt{results['Null']} & ndarray & $\min_k T_{\mu\nu}k^\mu k^\nu$;
                                      shape \texttt{(Nt,Nx,Ny,Nz)};
                                      negative = NEC violated \\
  \texttt{results['Weak']} & ndarray & $\min(\rho,\;\text{NEC})$;
                                      negative = WEC violated \\
  \texttt{T\_euler}        & ndarray & $\hat{T}^{\hat\mu\hat\nu}$ in
                                      Eulerian frame;
                                      \texttt{T\_euler[0,0]} = $\rho$ \\
  \bottomrule
\end{tabularx}
\end{center}

\paragraph{Null vector sampling.}
Test vectors are drawn from the Fibonacci lattice on $S^2$, giving a
near-uniform distribution.  For a reliable minimum, use $\geq 50$
vectors.  Increasing \texttt{num\_vecs} improves accuracy but increases
runtime linearly.

\subsection{Helper: \texttt{generate\_uniform\_field}}

\begin{lstlisting}
from warpfactory.analyzer.energy_conditions import generate_uniform_field

vecs = generate_uniform_field('nulllike', num_angular=100)
# vecs.shape == (4, 100)
# vecs[0, :] == 1.0   (time component)
# vecs[1:, :] are unit 3-vectors distributed on the sphere
\end{lstlisting}

\section{Frame Transformations}
\label{sec:api_transform}

\subsection{\texttt{do\_frame\_transfer}}

Transforms a stress-energy tensor to the Eulerian (locally flat) frame:

\begin{lstlisting}
from warpfactory.analyzer.transform import do_frame_transfer

energy_eu = do_frame_transfer(metric, energy_tensor, frame='Eulerian')
# energy_eu.index == 'contravariant'
# energy_eu.frame == 'Eulerian'
\end{lstlisting}

The transformation proceeds in three steps:
\begin{enumerate}
  \item Lower $T^{\mu\nu}$ to $T_{\mu\nu}$ using the curved metric.
  \item Compute the Cholesky-like tetrad matrix $M$ such that
        $M^T g M = \eta$.
  \item Apply: $\hat{T}_{ab} = M^\mu{}_a\,T_{\mu\nu}\,M^\nu{}_b$.
  \item Raise indices using the flat Minkowski metric: flip sign of
        $\hat{T}^{0i}$ and $\hat{T}^{i0}$.
\end{enumerate}

\subsection{\texttt{get\_eulerian\_transformation\_matrix}}

\begin{lstlisting}
from warpfactory.analyzer.transform import (
    get_eulerian_transformation_matrix
)

M = get_eulerian_transformation_matrix(metric.tensor)
# M.shape == (4, 4, Nt, Nx, Ny, Nz)
# Satisfies M^T g M = eta at every grid point
\end{lstlisting}

This function computes the Cholesky factorisation of the metric
component-by-component.  Each matrix element is expressed in closed
form using determinant cofactors.

\subsection{\texttt{change\_tensor\_index}}

\begin{lstlisting}
from warpfactory.analyzer.transform import change_tensor_index

T_cov = change_tensor_index(T_con, 'covariant', metric)
T_con = change_tensor_index(T_cov, 'contravariant', metric)
\end{lstlisting}

Raises or lowers both indices of a rank-2 tensor using the formula:
\begin{equation}
  T^{\mu\nu} = g^{\mu\alpha}\,T_{\alpha\beta}\,g^{\nu\beta}
\end{equation}
implemented as a batched matrix product over the grid axes.

\section{ADM Scalars}
\label{sec:scalars}

\subsection{\texttt{three\_plus\_one\_decomposer}}

\begin{lstlisting}
from warpfactory.analyzer.scalars import three_plus_one_decomposer

alpha, beta_down, gamma, beta_up, gamma_up = (
    three_plus_one_decomposer(metric)
)
# alpha.shape    == (Nt, Nx, Ny, Nz)   -- lapse
# beta_down.shape == (3, Nt, Nx, Ny, Nz) -- covariant shift
# gamma.shape    == (3,3,Nt,Nx,Ny,Nz)   -- spatial metric
\end{lstlisting}

Extracts 3+1 ADM variables from the full 4-metric:
\begin{align}
  \alpha &= \sqrt{\beta_i\beta^i - g_{00}}, &
  \beta_i &= g_{0i}, &
  \gamma_{ij} &= g_{ij}
\end{align}

\subsection{\texttt{calculate\_scalars}}

\begin{lstlisting}
from warpfactory.analyzer.scalars import calculate_scalars

scalars = calculate_scalars(metric, grid_scale)
# scalars['expansion'].shape == (Nt, Nx, Ny, Nz)
# scalars['shear'].shape     == (Nt, Nx, Ny, Nz)
# scalars['vorticity'].shape == (Nt, Nx, Ny, Nz)
\end{lstlisting}

Computes the kinematic quantities of the Eulerian congruence:
\begin{align}
  \theta &= \nabla_\mu n^\mu \quad (\text{expansion}) \\
  \sigma_{ab} &= \nabla_{(a}n_{b)} - \tfrac{1}{3}\theta h_{ab}
                \quad (\text{shear}) \\
  \omega_{ab} &= \nabla_{[a}n_{b]} \quad (\text{vorticity})
\end{align}
where $n^\mu = (1/\alpha, -\beta^i/\alpha)$ is the Eulerian 4-velocity
and $h_{ab}$ is the spatial projector.

\section{Momentum Flow Lines}

\begin{lstlisting}
from warpfactory.analyzer.momentum_flow import get_momentum_flow_lines

flow_lines = get_momentum_flow_lines(
    energy_tensor = T_euler[:, :, t_sl, :, :, :],  # 3D time slice
    start_points  = [(12, 12, 12), (15, 15, 15)],
    step_size     = 0.5,
    max_steps     = 200,
    scale_factor  = 1.0
)
# Returns list of (N_steps, 3) arrays — one streamline per start point
\end{lstlisting}

Integrates the momentum density $T^{0i}$ (energy flux) as a vector
field using trilinear interpolation at each step.

%% ════════════════════════════════════════════════════════════
\chapter{ADM Decomposition}
\label{chap:adm}
%% ════════════════════════════════════════════════════════════

\section{\texttt{decompose\_3plus1}}

\begin{lstlisting}
from warpfactory.algo.adm import decompose_3plus1

alpha, beta_list, gamma = decompose_3plus1(metric)
# alpha      : (Nt, Nx, Ny, Nz)
# beta_list  : [beta_x, beta_y, beta_z] each (Nt, Nx, Ny, Nz)
# gamma      : (3, 3, Nt, Nx, Ny, Nz)
\end{lstlisting}

Direct extraction from $g_{\mu\nu}$:
\begin{align*}
  \alpha &= \sqrt{-g_{00} + \gamma^{ij}\beta_i\beta_j}, &
  \beta_i &= g_{0i}, &
  \gamma_{ij} &= g_{ij}
\end{align*}

\section{\texttt{reconstruct\_3plus1}}

\begin{lstlisting}
from warpfactory.algo.adm import reconstruct_3plus1

g = reconstruct_3plus1(alpha, beta_list, gamma)
# Returns: (4, 4, Nt, Nx, Ny, Nz) metric tensor
\end{lstlisting}

Reconstructs the full 4-metric from ADM variables using
\cref{eq:adm_comp}.

%% ════════════════════════════════════════════════════════════
\chapter{Visualisation}
\label{chap:vis}
%% ════════════════════════════════════════════════════════════

\section{\texttt{plot\_tensor}}

\begin{lstlisting}
from warpfactory.visualizer.plotter import plot_tensor

plot_tensor(
    tensor          = metric,        # Metric object or raw ndarray
    slice_planes    = (1, 2),        # keep axes 1,2 (x-y plane)
    slice_locations = [3, None, None, None, 15],
                     # [t_idx, -, -, -, z_idx]; None = centre
    save_dir        = './plots/',    # None = display interactively
    filename_prefix = 'alcubierre'
)
\end{lstlisting}

Generates a $4\times4$ grid of subplots, one for each component
$g_{\mu\nu}$, showing a 2-D slice.  Upper-triangle components are
duplicated (exploiting symmetry).

\section{\texttt{plot\_3plus1}}

\begin{lstlisting}
from warpfactory.visualizer.plotter import plot_3plus1

plot_3plus1(
    metric          = metric,
    slice_planes    = (1, 2),
    slice_locations = None,          # all centre slices
    save_dir        = './plots/'
)
\end{lstlisting}

Produces three figure sets:
\begin{enumerate}
  \item Lapse function $\alpha$ (single panel).
  \item Shift vector $\beta^i$ (three-panel row for $x$, $y$, $z$).
  \item Spatial metric $\gamma_{ij}$ (six panels for unique components).
\end{enumerate}

\section{\texttt{get\_slice\_data}}

\begin{lstlisting}
from warpfactory.visualizer.plotter import get_slice_data

slice2d = get_slice_data(
    data            = metric.tensor[0, 0],  # (Nt, Nx, Ny, Nz)
    slice_planes    = (1, 2),               # keep x, y
    slice_locations = [3, None, None, 15]   # fix t=3, z=15
)
# Returns: (Nx, Ny) 2-D array
\end{lstlisting}

%% ════════════════════════════════════════════════════════════
\chapter{Google Colab Notebook}
\label{chap:colab}
%% ════════════════════════════════════════════════════════════

\section{Overview}

The file \texttt{WarpFactory\_Colab.ipynb} at the project root is a
fully self-contained Jupyter notebook designed for Google Colab.  All
physics code is inlined (no local package installation required), GPU
acceleration is handled automatically via CuPy, and code units
($C=G=1$) are used throughout.

\section{Quick Start on Colab}

\begin{enumerate}
  \item Upload \texttt{WarpFactory\_Colab.ipynb} to Google Drive or
        open it directly on GitHub via the \emph{Open in Colab} badge.
  \item Set \textbf{Runtime $\to$ Change runtime type $\to$ Hardware
        accelerator $\to$ GPU} (T4 or A100).
  \item Click \textbf{Runtime $\to$ Run all} (\texttt{Ctrl+F9}).
\end{enumerate}

\section{Notebook Structure}

\begin{center}
\begin{tabularx}{\linewidth}{rlX}
  \toprule
  Cell & Type & Content \\
  \midrule
  1 & Code & Install NumPy, matplotlib; detect GPU; install CuPy \\
  2 & Code & Core classes (\texttt{Metric}, constants) \\
  3 & Code & Grid utilities, shape functions \\
  4 & Code & Metric generators (Alcubierre, Schwarzschild,
             Van Den Broeck) \\
  5 & Code & Finite-difference solver \\
  6 & Code & Tensor calculations (Christoffel through $T^{\mu\nu}$) \\
  7 & Code & Frame transforms and energy conditions \\
  8 & Code & Demo 1 — Alcubierre warp drive \\
  9 & Code & Demo 1 — Visualisation (energy density + conditions) \\
  10 & Code & Demo 1 — Shape function plot \\
  11 & Code & Demo 2 — Schwarzschild black hole \\
  12 & Code & Demo 3 — Van Den Broeck metric \\
  13 & Code & GPU performance benchmark \\
  14 & Markdown & Tips, physics reference, further reading \\
  \bottomrule
\end{tabularx}
\end{center}

\section{GPU Support}

The notebook uses a global \texttt{xp} variable that points to either
CuPy (GPU) or NumPy (CPU):

\begin{lstlisting}[caption={GPU/CPU selection in notebook}]
try:
    import cupy as xp
    GPU_AVAILABLE = True
except ImportError:
    import numpy as xp
    GPU_AVAILABLE = False
\end{lstlisting}

All array operations in the inlined physics functions use \texttt{xp},
making them transparently GPU-accelerated when CuPy is available.  The
helper \texttt{to\_np(arr)} transfers arrays back to NumPy for
matplotlib plotting.

\section{Expected Performance}

\begin{center}
\begin{tabular}{lccc}
  \toprule
  Grid size & CPU (NumPy) & GPU (CuPy T4) & Speedup \\
  \midrule
  $20^3$  & $\sim$10 s  & $\sim$8 s  & $\sim$1.2$\times$ \\
  $40^3$  & $\sim$80 s  & $\sim$20 s & $\sim$4$\times$ \\
  $60^3$  & $\sim$5 min & $\sim$40 s & $\sim$8$\times$ \\
  \bottomrule
\end{tabular}
\end{center}

For grids smaller than $\approx 30^3$ the GPU overhead dominates and
CPU is competitive; for $50^3$ and above, GPU is strongly recommended.

%% ════════════════════════════════════════════════════════════
\chapter{Complete API Reference}
\label{chap:api}
%% ════════════════════════════════════════════════════════════

This chapter provides a consolidated API table for all public functions.

\begin{longtable}{p{4.5cm} p{3.5cm} p{6.5cm}}
  \toprule
  \textbf{Function / Class} &
  \textbf{Module} &
  \textbf{Signature / Return} \\
  \midrule
  \endfirsthead
  \multicolumn{3}{c}{\small\itshape (continued)} \\
  \toprule
  \textbf{Function / Class} &
  \textbf{Module} &
  \textbf{Signature / Return} \\
  \midrule
  \endhead
  \bottomrule
  \endfoot

  \texttt{Metric}
  & \texttt{generator.base}
  & Dataclass: \texttt{tensor, coords, scaling, name, index, params} \\[0.3em]

  \texttt{create\_grid}
  & \texttt{generator.commons}
  & \texttt{(grid\_size, grid\_scale) $\to$ dict} \\[0.3em]

  \texttt{get\_minkowski\_metric}
  & \texttt{generator.commons}
  & \texttt{(grid\_size) $\to$ ndarray(4,4,...)} \\[0.3em]

  \texttt{shape\_function\_alcubierre}
  & \texttt{generator.commons}
  & \texttt{(r, R, sigma) $\to$ ndarray} \\[0.3em]

  \texttt{create\_minkowski\_metric}
  & \texttt{generator.minkowski}
  & \texttt{(grid\_size, grid\_scale) $\to$ Metric} \\[0.3em]

  \texttt{create\_alcubierre\_metric}
  & \texttt{generator.alcubierre}
  & \texttt{(grid\_size, grid\_scale, world\_center, v, R, sigma) $\to$ Metric} \\[0.3em]

  \texttt{create\_schwarzschild\_metric}
  & \texttt{generator.schwarzschild}
  & \texttt{(grid\_size, grid\_scale, world\_center, rs) $\to$ Metric} \\[0.3em]

  \texttt{create\_van\_den\_broeck\_metric}
  & \texttt{generator.van\_den\_broeck}
  & \texttt{(\ldots, v, R1, $\sigma_1$, R2, $\sigma_2$, A) $\to$ Metric} \\[0.3em]

  \texttt{create\_lentz\_metric}
  & \texttt{generator.lentz}
  & \texttt{(grid\_size, grid\_scale, world\_center, v, scale) $\to$ Metric} \\[0.3em]

  \texttt{create\_warp\_shell\_metric}
  & \texttt{generator.warp\_shell}
  & \texttt{(\ldots, mass, r\_inner, r\_outer, r\_buff, sigma, \ldots) $\to$ Metric} \\[0.3em]

  \texttt{create\_modified\_time\_metric}
  & \texttt{generator.modified\_time}
  & \texttt{(\ldots, v, R, sigma, A) $\to$ Metric} \\[0.3em]

  \texttt{solve\_energy\_tensor}
  & \texttt{solver.solvers}
  & \texttt{(metric) $\to$ Metric[T\textsuperscript{$\mu\nu$}]} \\[0.3em]

  \texttt{verify\_tensor}
  & \texttt{solver.tensor\_utils}
  & \texttt{(tensor, raise\_error) $\to$ bool} \\[0.3em]

  \texttt{get\_c4\_inv}
  & \texttt{solver.tensor\_utils}
  & \texttt{(g) $\to$ ndarray g\textsuperscript{$\mu\nu$}} \\[0.3em]

  \texttt{get\_c3\_inv}
  & \texttt{solver.tensor\_utils}
  & \texttt{($\gamma$) $\to$ ndarray $\gamma^{ij}$} \\[0.3em]

  \texttt{get\_christoffel}
  & \texttt{solver.tensor\_utils}
  & \texttt{(g, grid\_scale) $\to$ ndarray(4,4,4,...)} \\[0.3em]

  \texttt{get\_ricci\_tensor}
  & \texttt{solver.tensor\_utils}
  & \texttt{(g, grid\_scale) $\to$ ndarray(4,4,...)} \\[0.3em]

  \texttt{get\_ricci\_scalar}
  & \texttt{solver.tensor\_utils}
  & \texttt{(R\_mn, g\_inv) $\to$ ndarray(...)} \\[0.3em]

  \texttt{get\_einstein\_tensor}
  & \texttt{solver.tensor\_utils}
  & \texttt{(R\_mn, R, g) $\to$ ndarray(4,4,...)} \\[0.3em]

  \texttt{get\_energy\_tensor}
  & \texttt{solver.tensor\_utils}
  & \texttt{(G\_mn, g\_inv) $\to$ ndarray(4,4,...)} \\[0.3em]

  \texttt{derive\_1st\_4th\_order}
  & \texttt{solver.finite\_difference}
  & \texttt{(A, axis, delta) $\to$ ndarray} \\[0.3em]

  \texttt{derive\_2nd\_4th\_order}
  & \texttt{solver.finite\_difference}
  & \texttt{(A, axis1, axis2, delta1[, delta2]) $\to$ ndarray} \\[0.3em]

  \texttt{calculate\_energy\_conditions}
  & \texttt{analyzer.energy\_conditions}
  & \texttt{(g, grid\_scale[, num\_vecs]) $\to$ (dict, ndarray)} \\[0.3em]

  \texttt{generate\_uniform\_field}
  & \texttt{analyzer.energy\_conditions}
  & \texttt{(field\_type, num\_angular) $\to$ ndarray(4,N)} \\[0.3em]

  \texttt{do\_frame\_transfer}
  & \texttt{analyzer.transform}
  & \texttt{(metric, energy\_tensor, frame) $\to$ Metric} \\[0.3em]

  \texttt{get\_eulerian\_transformation\_matrix}
  & \texttt{analyzer.transform}
  & \texttt{(g) $\to$ ndarray M} \\[0.3em]

  \texttt{change\_tensor\_index}
  & \texttt{analyzer.transform}
  & \texttt{(tensor, new\_index, metric) $\to$ Metric} \\[0.3em]

  \texttt{three\_plus\_one\_decomposer}
  & \texttt{analyzer.scalars}
  & \texttt{(metric) $\to$ ($\alpha$, $\beta_\downarrow$, $\gamma$, $\beta^\uparrow$, $\gamma^{ij}$)} \\[0.3em]

  \texttt{calculate\_scalars}
  & \texttt{analyzer.scalars}
  & \texttt{(metric, grid\_scale) $\to$ dict\{expansion, shear, vorticity\}} \\[0.3em]

  \texttt{get\_momentum\_flow\_lines}
  & \texttt{analyzer.momentum\_flow}
  & \texttt{(energy\_tensor, start\_pts, step, max\_steps, scale) $\to$ list} \\[0.3em]

  \texttt{plot\_tensor}
  & \texttt{visualizer.plotter}
  & \texttt{(tensor[, slice\_planes, locs, save\_dir, prefix])} \\[0.3em]

  \texttt{plot\_3plus1}
  & \texttt{visualizer.plotter}
  & \texttt{(metric[, slice\_planes, locs, save\_dir, prefix])} \\[0.3em]

  \texttt{get\_slice\_data}
  & \texttt{visualizer.plotter}
  & \texttt{(data, slice\_planes, locs) $\to$ ndarray 2D} \\[0.3em]

  \texttt{decompose\_3plus1}
  & \texttt{algo.adm}
  & \texttt{(metric) $\to$ ($\alpha$, [$\beta_i$], $\gamma$)} \\[0.3em]

  \texttt{reconstruct\_3plus1}
  & \texttt{algo.adm}
  & \texttt{($\alpha$, $\beta$, $\gamma$) $\to$ ndarray(4,4,...)} \\[0.3em]

  \texttt{tov\_const\_density}
  & \texttt{utils.helpers}
  & \texttt{(R, M, rho, r) $\to$ ndarray P} \\[0.3em]

  \texttt{alpha\_numeric\_solver}
  & \texttt{utils.helpers}
  & \texttt{(M, P, R, r) $\to$ ndarray $\alpha$} \\[0.3em]

  \texttt{compact\_sigmoid}
  & \texttt{utils.helpers}
  & \texttt{(r, R1, R2, sigma, Rbuff) $\to$ ndarray} \\[0.3em]

  \texttt{sph2cart\_diag}
  & \texttt{utils.helpers}
  & \texttt{($\theta$, $\phi$, g11, g22) $\to$ 7 ndarrays} \\[0.3em]

\end{longtable}

%% ════════════════════════════════════════════════════════════
\chapter{Practical Guidance}
\label{chap:practical}
%% ════════════════════════════════════════════════════════════

\section{Choosing Grid Parameters}

\begin{center}
\begin{tabularx}{\linewidth}{lX}
  \toprule
  Parameter & Recommendation \\
  \midrule
  $N_t$ & Minimum 5; use 7 for safe finite-difference boundary
           handling.  Static metrics can use $N_t = 1$ (no
           time derivatives needed). \\
  $N_x = N_y = N_z$ & $\ge 25$ for exploratory work; $\ge 50$ for
           publication-quality results.  Increase for sharper
           shells (large $\sigma$). \\
  $dx$ & Choose so $R \gg dx$ (at least 5 cells per bubble radius).
         Typical: $dx = R / 7$. \\
  $\sigma$ & Shell thickness $\sim R / \sigma$.  For $\sigma = 8$,
             $R=1.5$, the shell spans $\approx 0.2$ length units.
             Ensure $dx \ll 0.2$. \\
  \bottomrule
\end{tabularx}
\end{center}

\section{Unit Conventions}

\begin{infobox}
\textbf{Always use code units ($C=G=1$) for exploratory simulations.}

In SI units, $C = 2.998 \times 10^8$ m/s.  Setting $v = 0.5$ with SI
constants moves the bubble $1.5 \times 10^8$ m per second, exiting a
6-metre grid in $40$ nanoseconds ($\Delta t < 10^{-7}$ required for
numerical stability).  Code units avoid this entirely.
\end{infobox}

To use code units: edit \texttt{warpfactory/constants.py} or, in the
Colab notebook, the constants cell sets \texttt{C = G = 1.0} directly.

\section{Interpreting Energy Condition Results}

\begin{center}
\begin{tabularx}{\linewidth}{lcX}
  \toprule
  Quantity & Sign & Interpretation \\
  \midrule
  \texttt{results['Null']} & $>0$ & NEC satisfied at this point \\
  \texttt{results['Null']} & $<0$ & NEC \emph{violated}; exotic matter needed \\
  \texttt{results['Weak']} & $<0$ & WEC violated; negative energy density \\
  \texttt{T\_euler[0,0]}  & $<0$ & Negative energy density $\rho < 0$ \\
  \bottomrule
\end{tabularx}
\end{center}

\paragraph{Expected results by metric:}
\begin{itemize}
  \item \textbf{Minkowski}: all quantities $\approx 0$.
  \item \textbf{Alcubierre}: large NEC and WEC violations in the bubble
        shell, concentrated in a ring around the $x$-axis.
  \item \textbf{Van Den Broeck}: similar violations but significantly
        smaller magnitude due to spatial expansion.
  \item \textbf{Schwarzschild}: all quantities $\approx 0$ (vacuum).
\end{itemize}

\section{Memory Usage}

The main memory bottleneck is the Christoffel symbol array
$\varGamma^\lambda_{\mu\nu}$, which has shape
$(4, 4, 4, N_t, N_x, N_y, N_z)$.  For float64:

\begin{center}
\begin{tabular}{lccc}
  \toprule
  Grid & $\varGamma$ size & $R_{\mu\nu}$ size & Total (approx.) \\
  \midrule
  $7 \times 25^3$ & 55 MB  & 5 MB  & $\sim$200 MB \\
  $7 \times 50^3$ & 440 MB & 40 MB & $\sim$2 GB \\
  $7 \times 75^3$ & 1.5 GB & 130 MB& $\sim$7 GB \\
  \bottomrule
\end{tabular}
\end{center}

Colab's T4 GPU has 16 GB VRAM — sufficient for $\sim 60^3$ grids.

\section{Common Errors}

\begin{center}
\begin{tabularx}{\linewidth}{lX}
  \toprule
  Error & Cause and Fix \\
  \midrule
  \texttt{ValueError: Tensor must be (4,4,...)} &
  Wrong array shape passed to a solver function.
  Ensure \texttt{metric.tensor.shape[:2] == (4,4)}. \\[0.3em]

  All-zero energy tensor &
  Grid too small ($N < 5$ in some dimension).
  The finite-difference function returns zeros silently for small arrays. \\[0.3em]

  \texttt{LinAlgError: Singular matrix} &
  Metric is degenerate (e.g.\ near horizon or at $r=0$).
  Add an epsilon guard or move the grid away from the singularity. \\[0.3em]

  NaN / Inf in results &
  Metric ill-conditioned (sharp discontinuity or very large $\sigma$).
  Increase $N$ or reduce $\sigma$. \\[0.3em]

  Slow computation &
  Large grid on CPU. Use GPU (\texttt{cupy}) or reduce $N$. \\
  \bottomrule
\end{tabularx}
\end{center}

%% ════════════════════════════════════════════════════════════
\chapter{Testing and Validation}
\label{chap:tests}
%% ════════════════════════════════════════════════════════════

All validation scripts are in the \texttt{tests/} directory.

\begin{center}
\begin{tabularx}{\linewidth}{llX}
  \toprule
  Script & Metric & Validates \\
  \midrule
  \texttt{verify\_package.py}         & Minkowski
  & $T^{\mu\nu} \approx 0$; smoke test for correct installation \\
  \texttt{validate\_alcubierre.py}    & Alcubierre
  & Negative energy density detected; NEC/WEC maps generated \\
  \texttt{validate\_schwarzschild.py} & Schwarzschild
  & Ricci scalar $R \approx 0$ (vacuum solution) \\
  \texttt{validate\_energy\_conditions.py} & Alcubierre
  & NEC violations $<-10^{-10}$ \\
  \texttt{validate\_solver.py}        & Alcubierre
  & Solver returns \texttt{Stress-Energy} type with finite values \\
  \texttt{validate\_lentz.py}         & Lentz
  & Metric creation and solver run without error \\
  \texttt{validate\_van\_den\_broeck.py} & Van Den Broeck
  & Metric creation successful \\
  \texttt{validate\_adm.py}           & Alcubierre
  & ADM decompose/reconstruct round-trips successfully \\
  \texttt{validate\_visualizer.py}    & Any
  & Plots saved without error \\
  \bottomrule
\end{tabularx}
\end{center}

\paragraph{Running all tests:}
\begin{lstlisting}[language=bash]
cd pyWarpFactory
python tests/verify_package.py
python tests/validate_alcubierre.py
python tests/validate_schwarzschild.py
python tests/validate_energy_conditions.py
\end{lstlisting}

%% ════════════════════════════════════════════════════════════
\chapter{Complete Worked Examples}
\label{chap:examples}
%% ════════════════════════════════════════════════════════════

\section{Example 1: Full Alcubierre Analysis}

\begin{lstlisting}[caption={Complete Alcubierre analysis script}]
"""
Full Alcubierre warp drive analysis:
 - Generate metric
 - Solve Einstein equations
 - Evaluate energy conditions
 - ADM decomposition
 - Plot all results
"""
import numpy as np
import matplotlib.pyplot as plt

from warpfactory.generator.alcubierre import create_alcubierre_metric
from warpfactory.solver.solvers import solve_energy_tensor
from warpfactory.analyzer.energy_conditions import (
    calculate_energy_conditions
)
from warpfactory.algo.adm import decompose_3plus1
from warpfactory.visualizer.plotter import plot_tensor, plot_3plus1

# ── Parameters ────────────────────────────────────────────────
Nt, N = 7, 30
grid_size  = (Nt, N, N, N)
grid_scale = (1.0, 0.3, 0.3, 0.3)
world_center = (1.0, N*0.3/2, N*0.3/2, N*0.3/2)

v, R, sigma = 0.5, 1.5, 8.0

# ── Metric ────────────────────────────────────────────────────
print("Generating Alcubierre metric …")
metric = create_alcubierre_metric(
    grid_size, grid_scale, world_center, v, R, sigma
)
print(f"  Metric shape: {metric.tensor.shape}")

# ── Einstein solver ───────────────────────────────────────────
print("Solving Einstein field equations …")
energy = solve_energy_tensor(metric)
print(f"  Max |T^00|: {np.max(np.abs(energy.tensor[0,0])):.4e}")

# ── Energy conditions ─────────────────────────────────────────
print("Evaluating energy conditions …")
results, T_euler = calculate_energy_conditions(
    metric.tensor, metric.scaling
)
t_sl, z_sl = Nt // 2, N // 2
rho = T_euler[0, 0, t_sl, :, :, z_sl]
nec = results['Null'][t_sl, :, :, z_sl]
wec = results['Weak'][t_sl, :, :, z_sl]
print(f"  rho min = {rho.min():.4e}")
print(f"  NEC min = {nec.min():.4e}")
print(f"  WEC min = {wec.min():.4e}")

# ── ADM decomposition ─────────────────────────────────────────
print("ADM decomposition …")
alpha, beta_list, gamma = decompose_3plus1(metric)
print(f"  Lapse alpha range: [{alpha.min():.4f}, {alpha.max():.4f}]")

# ── Plot ──────────────────────────────────────────────────────
fig, axes = plt.subplots(2, 3, figsize=(16, 10))
ext = [0, N*0.3, 0, N*0.3]

axes[0,0].imshow(rho.T, origin='lower', extent=ext, cmap='RdBu_r')
axes[0,0].set_title(r'Energy density $\rho = \hat{T}^{00}$')

axes[0,1].imshow(nec.T, origin='lower', extent=ext, cmap='RdBu_r')
axes[0,1].set_title('NEC violation map')

axes[0,2].imshow(wec.T, origin='lower', extent=ext, cmap='RdBu_r')
axes[0,2].set_title('WEC violation map')

shift = beta_list[0][t_sl, :, :, z_sl]
axes[1,0].imshow(shift.T, origin='lower', extent=ext, cmap='coolwarm')
axes[1,0].set_title(r'Shift $\beta_x$')

alp_sl = alpha[t_sl, :, :, z_sl]
axes[1,1].imshow(alp_sl.T, origin='lower', extent=ext, cmap='viridis')
axes[1,1].set_title(r'Lapse $\alpha$')

gxx = gamma[0, 0, t_sl, :, :, z_sl]
axes[1,2].imshow(gxx.T, origin='lower', extent=ext, cmap='plasma')
axes[1,2].set_title(r'Spatial metric $\gamma_{xx}$')

for ax in axes.flat:
    ax.set_xlabel('x'); ax.set_ylabel('y')

plt.suptitle(f'Alcubierre  v={v}  R={R}  sigma={sigma}',
             fontsize=13, fontweight='bold')
plt.tight_layout()
plt.savefig('alcubierre_full_analysis.png', dpi=130)
plt.show()
print("Saved: alcubierre_full_analysis.png")
\end{lstlisting}

\section{Example 2: Comparing Alcubierre and Van Den Broeck}

\begin{lstlisting}[caption={Side-by-side energy comparison}]
import numpy as np
from warpfactory.generator.alcubierre import create_alcubierre_metric
from warpfactory.generator.van_den_broeck import (
    create_van_den_broeck_metric
)
from warpfactory.analyzer.energy_conditions import (
    calculate_energy_conditions
)
import matplotlib.pyplot as plt

N = 30
grid_size    = (7, N, N, N)
grid_scale   = (1.0, 0.3, 0.3, 0.3)
world_center = (1.0, 4.5, 4.5, 4.5)
t_sl, z_sl   = 3, N // 2

# Alcubierre
m_alc = create_alcubierre_metric(
    grid_size, grid_scale, world_center, 0.5, 1.5, 8.0
)
_, T_alc = calculate_energy_conditions(m_alc.tensor, grid_scale)

# Van Den Broeck (same velocity, similar inner radius)
m_vdb = create_van_den_broeck_metric(
    grid_size, grid_scale, world_center,
    v=0.5, R1=2.5, sigma1=6.0, R2=1.0, sigma2=8.0, A=2.0
)
_, T_vdb = calculate_energy_conditions(m_vdb.tensor, grid_scale)

fig, axes = plt.subplots(1, 2, figsize=(13, 5))
ext = [0, N*0.3, 0, N*0.3]
vmax = max(abs(T_alc[0,0,t_sl,:,:,z_sl]).max(),
           abs(T_vdb[0,0,t_sl,:,:,z_sl]).max())

axes[0].imshow(T_alc[0,0,t_sl,:,:,z_sl].T, origin='lower',
               extent=ext, cmap='RdBu_r', vmin=-vmax, vmax=vmax)
axes[0].set_title('Alcubierre — energy density')

axes[1].imshow(T_vdb[0,0,t_sl,:,:,z_sl].T, origin='lower',
               extent=ext, cmap='RdBu_r', vmin=-vmax, vmax=vmax)
axes[1].set_title('Van Den Broeck — energy density')

plt.tight_layout()
plt.savefig('comparison.png', dpi=130)
print("Saved comparison.png")
\end{lstlisting}

%% ════════════════════════════════════════════════════════════
\appendix
%% ════════════════════════════════════════════════════════════

\chapter{Physical Constants}
\label{app:constants}

\begin{center}
\begin{tabular}{llll}
  \toprule
  Symbol & SI Value & Unit & Python name \\
  \midrule
  $c$ & $2.99792458 \times 10^8$ & m s$^{-1}$ & \texttt{C} \\
  $G$ & $6.67430   \times 10^{-11}$ & m$^3$ kg$^{-1}$ s$^{-2}$ & \texttt{G} \\
  \bottomrule
\end{tabular}
\end{center}

In code units: $C = G = 1$.  The stress-energy tensor scale factor is:
\begin{equation}
  \frac{c^4}{8\pi G}
  \approx 4.815 \times 10^{42} \quad \text{(SI: J m}^{-3}\text{)}
\end{equation}

\chapter{Tensor Index Notation}
\label{app:indices}

\begin{center}
\begin{tabular}{lcl}
  \toprule
  Python axis & Index & Physical meaning \\
  \midrule
  0 & $\mu = 0$ & time $t$ \\
  1 & $\mu = 1$ & spatial $x$ \\
  2 & $\mu = 2$ & spatial $y$ \\
  3 & $\mu = 3$ & spatial $z$ \\
  \midrule
  2 (grid dim) & — & $t$ grid \\
  3 (grid dim) & — & $x$ grid \\
  4 (grid dim) & — & $y$ grid \\
  5 (grid dim) & — & $z$ grid \\
  \bottomrule
\end{tabular}
\end{center}

\paragraph{Array layout:}
\texttt{metric.tensor[mu, nu, t\_idx, x\_idx, y\_idx, z\_idx]}

\chapter{Fourth-Order Stencil Derivation}
\label{app:fd}

The 4th-order central difference stencil for the first derivative is
derived from the Taylor expansion:
\begin{align*}
  f_{i+1} &= f_i + f'_i h + \tfrac{1}{2}f''_i h^2 + \cdots \\
  f_{i+2} &= f_i + 2f'_i h + 2f''_i h^2 + \cdots
\end{align*}
Eliminating higher derivatives yields:
\begin{equation}
  f'_i = \frac{-f_{i+2} + 8f_{i+1} - 8f_{i-1} + f_{i-2}}{12h}
  + \mathcal{O}(h^4)
\end{equation}
For the second derivative (unmixed):
\begin{equation}
  f''_i = \frac{-f_{i+2} + 16f_{i+1} - 30f_i + 16f_{i-1} - f_{i-2}}{12h^2}
  + \mathcal{O}(h^4)
\end{equation}

\chapter{Changelog}
\label{app:changelog}

\begin{center}
\begin{tabularx}{\linewidth}{llX}
  \toprule
  Version & Date & Changes \\
  \midrule
  1.0.0 & 2026-02-28 &
  Initial Python release: Alcubierre, Schwarzschild, Van Den Broeck,
  Lentz, Warp Shell, Modified Time metrics; full Einstein solver;
  NEC/WEC energy conditions; ADM decomposition; Colab notebook. \\[0.3em]

  1.0.1 & 2026-02-28 &
  Bug fixes: (i) removed double-count in \texttt{get\_ricci\_tensor}
  term-4 loop; (ii) rewrote \texttt{calculate\_energy\_conditions} to
  fix API mismatch and NameError in WEC block; (iii) cleaned up
  \texttt{validate\_alcubierre.py}; (iv) rebuilt Colab notebook with
  GPU support. \\
  \bottomrule
\end{tabularx}
\end{center}

\chapter{Further Reading}
\label{app:refs}

\begin{description}
  \item[Alcubierre (1994)]
    M.~Alcubierre, \textit{The warp drive: hyper-fast travel within
    general relativity}, Class.\ Quantum Grav.\ \textbf{11}, L73 (1994).

  \item[Van Den Broeck (1999)]
    C.~Van Den Broeck, \textit{A warp drive with more reasonable total
    energy requirements}, Class.\ Quantum Grav.\ \textbf{16}, 3973 (1999).

  \item[Lentz (2021)]
    E.~W.~Lentz, \textit{Breaking the warp barrier: hyper-fast solitons
    in Einstein-Maxwell-plasma theory}, Class.\ Quantum Grav.\ \textbf{38},
    075015 (2021).

  \item[Arnowitt, Deser, Misner (1962)]
    R.~Arnowitt, S.~Deser, C.~W.~Misner, \textit{The dynamics of
    general relativity}, Gravitation, ed.\ Witten, Wiley (1962).

  \item[Misner, Thorne, Wheeler (1973)]
    C.~W.~Misner, K.~S.~Thorne, J.~A.~Wheeler, \textit{Gravitation},
    Freeman (1973). — Standard GR reference.

  \item[Baumgarte \& Shapiro (2010)]
    T.~W.~Baumgarte, S.~L.~Shapiro, \textit{Numerical Relativity:
    Solving Einstein's Equations on the Computer}, Cambridge (2010).
    — Background on 3+1 formalism and finite differences.
\end{description}

\end{document}
